
Digital services often feel "weightless". They are something you use daily without thinking much about the cost of them. They are services where a User Experience Designer tried to provide the best possible experience for their users. In their work process they have to make multiple choices, which in reality are not weightless at all. For example, a Netflix designer setting the default video quality to the highest possible one with autoplay enabled makes a choice that will be replicated for more than 325 million paid subscribers streaming over 95 billion hours in the first half of 2025 \cite{Netflix2025Q4, Netflix2025H1}. A product manager deciding to implement an AI chatbot feature because it is a new trend, instead of a keyword search function, claims a share of data center infrastructure that consumes electricity from grids, still 58,9\% powered by fossil fuels in 2024 in the United States \cite{IEA2025GlobalReview}.

\par While in 2024 data centers accounted for approximately 1.5\% of global electricity consumption multiple factors make this ethically significant. Depending on different scenario cases this consumption share is predicted to rise to 2\% - 4.4\%, potentially tripling in a decade by 2035 \cite{IEA2025EnergyAI}. Additionally data centres consume fresh water. In their 2025 environmental report Google claimed their total data centre water consumption in the year 2024 to be 29.5 billion litres of water growing yearly \cite{Google2025EnvReport}. Lastly, this growth occurs against the backdrop of a climate crisis. In 2023 the Intergovernmental Panel on Climate Change warned that human-caused global warming has already reached 1.1 °C, triggering widespread and rapid changes across the biosphere that pose a direct threat to human well-being and planetary health \cite{IPCC2023Synthesis}. Another report shows that of 34 tracked planetary vital signs, 22 are at record levels, many showing alarming trends \cite{Ripple2025StateClimate}. The World Meteorological Organization confirmed that including 2025 the eleven past years were the 11 warmest years on record \cite{WMO2026WarmestYear}. In addition, a 2026 report from the United Nations Universities Institute for Water, Environment and Health argues that for parts of the world the water crisis should rather be called water bankruptcy, since a crisis would suggest the possibility of turning back to normalcy \cite{UNU2026WaterBankruptcy}. These converging environmental pressures demonstrate the urgency of climate responsibility in all sectors, including digital UX design.

\par This paper argues that UX designers bear professional responsibility for
environmental consequences that scale across billions of users, yet face barriers preventing ethical action. First, lack of transparency: tech
companies treat energy consumption data as proprietary information, making
informed design decisions impossible. Second, structural pressure: designers are evaluated on
engagement metrics (time-on-site, retention), not environmental impact.
Proposing to disable autoplay or default to lower video quality means
accepting reduced engagement scores and potential career consequences. Individual designers retain
agency to advocate for sustainable choices within their professional sphere, but
cannot overcome these barriers alone. The normative conclusion is therefore:
designers should try to minimize environmental impact at every decision point and support
regulatory interventions (mandatory energy disclosure, eco-friendly default
standards) that make sustainable design professionally viable rather than
professionally punished.
\par This paper proceeds in five parts. Section II establishes the materiality
problem: the digital space is not weightless. Different UX design decisions have energy impacts that scale across billions of users.
Section III analyzes this through ethical frameworks of climate justice and autonomy, demonstrating that current practices violate
these principles by externalizing harm and denying both users and designers
the information necessary for informed choice. Section IV examines the
designer's professional dilemma in depth, demonstrating that even ethically-informed
designers face severe structural constraints: business models reward
engagement maximization regardless of environmental cost.
Finally, Section V proposes normative recommendations.
